\scanpage{103}
\mspage{102}
Je vais prendre comme relations primitives dans $\unsure{\mathcal M|M}$
\[
\leq,\quad\ll,\quad\mathrel{|\circ|}
\]
Avec ça, on va définir $\unsure{\gtrless|\bowtie}$ ainsi (version \deleted{``\unsure{optimiste}''} ``magasin spacieux'', on reviendra bientôt sur la version ``magasin polyédral'') : pour $X,Y\in\mathcal M$, posant $L=\widetilde X\cap\widetilde Y$ \inserted{plus grande partie fermée commune}
\[
X\gtrless Y\quad\Longleftrightarrow\quad\forall X'\in\widetilde X\setminus L,\quad Y'\in\widetilde Y\setminus L,\quad X'\mathrel{|\circ|}Y'
\]
% slanted
\margin{cette définition ne fait pas intervenir $\ll$.}
Nous allons d'abord poser les axiomes Mag $\Lambda1\ldots\inserted{\Lambda\unsure{6|5}}$ (anciennement Prat $\Lambda1\unsure{\text{?}|\S}\ldots$) de la page 51, avec $\Lambda=\mathcal M$.

Mag $\mathcal M1$, Mag $\mathcal M3$, Mag $\mathcal M4$ \unsure{ne varient} \inserted{donnant Mag $1''$, $2''$, $3''$}

Mag $\mathcal M2$ tombe (car \inserted{$\unsure{\mathcal M\Lambda2|\mathrm{M}\Lambda2}$} tautologique si $\Lambda=\mathcal M$)

Mag $\mathcal M5$ \deleted{est tautologique} va résulter du nouvel axiome (remplaçant At L3)
\[
\text{\uline{Mag} }4''\qquad\forall X\in\mathcal M,\ \text{et }Y,Z\leq X,\ \inserted{Y\ne Z,}\ \text{on a}\quad Y\deleted{\unsure{\gtrless}}\mathrel{|\circ|}Z
\]

\diagram{\begin{tikzcd}[ampersand replacement=\&,column sep=small,row sep=small]
\&\&\widetilde X\arrow[dl,dash]\arrow[dr,dash]\&\&\widetilde Y\arrow[dl,dash]\arrow[dr,dash]\&\&\\
\&\widetilde X'\arrow[dr,dash]\arrow[dl,phantom,"\ni" description,sloped]\&\&L\arrow[dl,dash]\arrow[dr,dash]\&\&\widetilde Y'\arrow[dl,dash]\arrow[dr,phantom,"\ni" description,sloped]\&\\
X''\&\&X'_L\arrow[dr,dash]\&\&Y'_L\arrow[dl,dash]\&\&Y''\\
\&\&\&L'\&\&\&
\end{tikzcd}}
Supposons $X\gtrless Y$, soient $X'\unsure{\nleq|\leq}X$, $Y'\leq Y$, prouvons $X'\gtrless Y'$. Introduisons $\longrightarrow$
\[
\begin{aligned}
\deleted{\unsure{X}}L&=\widetilde X\cap\widetilde Y,\quad L'=\widetilde X'\cap\widetilde Y',\quad X'_L=\widetilde X'\cap L=\widetilde X'\cap\widetilde Y,\\
Y'_L&=\widetilde Y'\cap L=\widetilde Y'\cap\widetilde X,\quad\text{donc }L'=X'_L\cap Y'_L.
\end{aligned}
\]
Soit $X''\in\widetilde X'\setminus L'$, $Y''=\widetilde Y'\setminus L'$, prouvons $X''\mathrel{|\circ|}Y''$. Si $X'',Y''\notin L$, cela résulte de l'hypothèse que $\deleted{\unsure{Z}}X\gtrless Y$. Si $X''\in L$ i.e. $X''\in\deleted{\unsure{\widetilde X_L}}\unsure{\widetilde Y}$, ($Y''$ non \unsure{entrant|partant}) alors $X'',Y''\in\widetilde Y$ et $X''\ne Y''$, donc par Mag $4''$ on gagne encore.

Il faut encore \unsure{imposer} Mag $\mathcal M6$ (ex Prat $\Lambda\unsure{6|5}$, pour $\Lambda=\mathcal M$).

Cor. de \uline{Mag} $4''$ Si $X,Y$ avec $X\gtrless Y$, $X\ne Y$, alors $X\mathrel{|\circ|}Y$.
